\section{SF-E847: scalar complete-intersection point extensions}\label{SF-E847:section}

\textbf{Checked mathematical status, 2026-11-22.} The author, root,
and independent reviewer \texttt{/root/g34\_eightfolds} have read the
complete frozen proof. All mathematical reviews passed, with no repair.
The independent review was completed at \texttt{02:53:53 UTC}, and its
full report was saved at \texttt{02:55:19 UTC}. The frozen Markdown
original retains its author-review status. This conversion has a separate
whole-text verification record. No numerical certificate or strict
comparison with every earlier local criterion is claimed.
Author: \texttt{/root/sixfolds}, GPT-6 Astra Pro; inherited effort is not exposed.
This argument does not use a noncommutative deformation, a homogeneous
grading, a real-multiplication action, or a Chern-character identity.
The actual finite-action scope remains the translation/tensor framework
of supplied DEG\_L S21. A strict comparison with all earlier local criteria
is a separate question and is not asserted in this draft.

\subsection{1. Statement and the scalar local type}\label{SF-E847:s1}

Let \(R=\mathbf C[[x_1,x_2,x_3,x_4]]\), with maximal ideal \(\mathfrak m\). Let
\(I=(f_1,f_2,f_3,f_4)\) be a proper \(\mathfrak m\)-primary complete-intersection ideal and
\(A=R/I\), of positive length \(\ell\). No restriction is made on the orders,
leading forms, or nonlinear terms of the \(f_i\). The ordinary Koszul
resolution gives

\[
 \operatorname{Hom}_R(I,A)=A^4,\qquad
 \operatorname{Ext}^1_R(I,A)=A^6.
 \tag{SC.1}
\]

Fix any nonzero scalar six-tuple \(b\) in \(A^6\), with arbitrary coefficients,
and let \(M=M_b\) be its actual extension

\[
 0\longrightarrow A\longrightarrow M
       \longrightarrow I\longrightarrow0.
 \tag{SC.2}
\]

The adjective scalar means that both the torsion and hull quotient are
one copy of \(A\); \(b\) is not a matrix-valued extension class. The map
\(\phi:M\to I\to R\) is the canonical evaluation into the double dual. Indeed
\(\operatorname{Hom}_R(A,R)=\operatorname{Ext}^1_R(A,R)=0\), and the sequence \(I\to R\to A\) gives
\(\operatorname{Hom}_R(I,R)=R\). Thus \(\operatorname{Hom}_R(M,R)=R\), \(M^{**}=R\), and evaluation is exactly \(\phi\).
Its kernel and cokernel are both \(A\). In particular \(M\) is rank one,
has nonzero finite torsion, and is free away from the closed point.

\textbf{Local dimension theorem SF-E847a.} With

\[
 a=\dim_{\mathbf C}\operatorname{Ann}_A(b),\qquad
 k=\dim_{\mathbf C}\operatorname{Hom}_R(M,A),\qquad
 e_i=\dim_{\mathbf C}\operatorname{Ext}^i_R(M,M)\quad(i>0),
 \tag{SC.3}
\]

one has

\[
 k=4\ell+a,\qquad e_1\ge9\ell+a,\qquad
 e_2=2e_1-2k\ge10\ell.
 \tag{SC.4}
\]

The positive \(\operatorname{Ext}\) groups are finite-dimensional because \(M\) is free off \(\mathfrak m\).
The identity in (SC.4) is proved from a finite-supported hull cone.
It is not an ordinary local self-duality assertion for \(R\operatorname{Hom}(M,M)\),
whose degree-zero endomorphism module has rank one.

\textbf{Global theorem SF-E847b.} Let \(X\) be an abelian fourfold and \(E\) a coherent
sheaf whose canonical map to its double dual has locally free hull \(H\)
and finite kernel and cokernel. Suppose at one point \(p\) there is an
actual isomorphism of completed modules

\[
 \widehat E_p\simeq M_b^{\oplus m}\oplus R^{\oplus s},
 \qquad m>0,\quad s\ge0,
 \tag{SC.5}
\]

for the scalar type (SC.2). Then, for every specified effective coherent
finite translation/tensor action \(G\) in the supplied S21 framework, the
kernel of the full middle semiregularity map on invariant \(\operatorname{Ext}\) contains
a subspace of dimension at least \(10\ell-6\). In particular \(E\) is not weakly
\(G\)-semiregular. This statement requires no balanced character, RM action,
polarization, B-field, simplicity, or second-order lifting implication.

The balanced-RM application is simply the specialization to the objects
sought in the main seed problem. It introduces no further hypothesis
into the proof of the global theorem.

\subsection{2. The finite endomorphism correction}\label{SF-E847:s2}

Set \(K=\operatorname{Hom}_R(M,A)\). Apply \(\operatorname{Hom}_R(-,A)\) to (SC.2). The beginning of the
long exact sequence is

\[
 0\longrightarrow A^4\longrightarrow K\longrightarrow A
        \xrightarrow{\ t\mapsto tb\ } A^6.
 \tag{SC.6}
\]

This proves \(k=4\ell+a\). Every endomorphism of \(M\) induces an endomorphism
of its canonical double dual \(R\). The resulting map \(\operatorname{End}_R(M)\to R\) is
surjective: multiplication by any \(r\) in \(R\) on \(M\) maps to multiplication
by \(r\) on \(R\). Its kernel consists exactly of the maps \(M\to A\to M\). Consequently

\[
 \operatorname{End}_R(M)\simeq R\oplus K
       \quad\hbox{as }R\hbox{-modules}.
 \tag{SC.7}
\]

The splitting is the scalar-multiplication map. It is crucial here
that the image in \(R\) is all of \(R\), not an unspecified ideal.
No finite-dimensional Euler characteristic is assigned to \(\operatorname{End}_R(M)\).

We will also need that \(\operatorname{End}_R(M)\) is local. An endomorphism induces
multiplication by \(r\) in \(R\) on the hull and by \(t\) in \(A\) on the torsion.
Compatibility with the actual extension class gives \((t-r\bmod I)b=0\).
Because \(b\) is nonzero, \(t-r\bmod I\) cannot be a unit. Their residues agree.
If the common residue is nonzero, the torsion and hull-image maps
are both invertible, so the endomorphism of \(M\) is invertible by (SC.2).
If the residue is zero, its difference from the identity is invertible.
This proves locality, with residue field \(\mathbf C\). Thus \(M\) is indecomposable
and is not isomorphic to \(R\).

\subsection{3. A relative Koszul family and its canonical hull}\label{SF-E847:s3}

Let \(D=\mathbf C[\varepsilon]/(\varepsilon^2)\) and \(S=D[[x_1,x_2,x_3,x_4]]\). Choose fixed
\(\mathbf C\)-linear lifts of \(A\) to \(R\) for the following finite-dimensional
parameters. For \(\alpha=(\alpha_i)\) in \(A^4\) and \(\beta=(\beta_{ij})\) in \(A^6\) set

\[
 f_i^{\epsilon}=f_i+\epsilon\widetilde\alpha_i,\qquad
 I_\epsilon=(f_1^\epsilon,\ldots,f_4^\epsilon),\qquad
 A_\epsilon=S/I_\epsilon,
 \qquad
 b^\epsilon_{ij}=\widetilde b_{ij}
                 +\epsilon\widetilde\beta_{ij}\pmod{I_\epsilon}.
 \tag{SC.8}
\]

All these are deformations over the fixed commutative base algebra \(S\).
No coordinate change of \(R\) is included in the equivalence relation.

Here is an elementary flatness check. The Koszul complex of \(f^\varepsilon\)
has a short exact sequence of complexes whose first map is multiplication
by \(\varepsilon\) and whose quotient is reduction modulo \(\varepsilon\):

\[
 0\longrightarrow K(f;R)\longrightarrow K(f^\epsilon;S)
       \longrightarrow K(f;R)\longrightarrow0.
 \tag{SC.9}
\]

The original Koszul complex has cohomology only in degree zero, equal
to \(A\). Its long exact cohomology sequence proves that the middle complex
is also acyclic in negative degrees, with a short exact sequence
\(0\to A\to A_\varepsilon\to A\to0\). Multiplication by \(\varepsilon\) on \(A_\varepsilon\) is the
composite of the last quotient with the first inclusion. Thus
\(\ker(\varepsilon)=\varepsilon A_\varepsilon\). Over \(D\) this is precisely flatness:
choose a \(\mathbf C\)-basis of \(\varepsilon A_\varepsilon\) and arbitrary lifts; these form
a \(D\)-basis. In particular \(A_\varepsilon\) is finite free of rank \(\ell\) over \(D\).
The kernel \(I_\varepsilon\) of \(S\to A_\varepsilon\) is \(D\)-flat as well, by the exact
sequence of these two flat modules.

The self-dual Koszul resolution just proved gives

\[
 \operatorname{Ext}^j_S(A_\epsilon,S)=0\quad(0\le j<4),
 \qquad
 \operatorname{Hom}_S(I_\epsilon,S)=S.
 \tag{SC.10}
\]

The latter isomorphism is restriction of scalar multiplication on \(S\).
The truncated Koszul resolution gives \(\operatorname{Ext}^1_S(I_\varepsilon,A_\varepsilon)\)=
\(A_\varepsilon^6\), so \(b^\varepsilon\) defines an actual extension
\(0\to A_\varepsilon\to M_\varepsilon\to I_\varepsilon\to0\). Its end terms are \(D\)-flat;
therefore \(M_\varepsilon\) is \(D\)-flat and reduces to the fixed extension \(M\).

For clarity this family can be constructed by free resolutions, without
a moduli-space smoothness assertion. Let \(d_\varepsilon\) be the Koszul
differential and let \(B_\varepsilon\) be contraction by the even two-form
\(b^\varepsilon\). The differential on the triangular resolution is
\(\left[\begin{smallmatrix}d_\varepsilon&B_\varepsilon\\0&-d_\varepsilon\end{smallmatrix}\right]\), with the same exterior modules
as the resolution of \(M\). Even contraction commutes with \(d_\varepsilon\),
so its square is zero. The short exact sequence with the two Koszul
complexes proves its exactness and the indicated extension in degree zero.
The coefficients of the first-order change are linear in \(\alpha\) and \(\beta\).

There is no unproved base-change assertion about double duals. Applying
\(\operatorname{Hom}_S(-,S)\) to the constructed extension and using (SC.10) gives,
canonically through its quotient map,

\[
 \operatorname{Hom}_S(M_\epsilon,S)
    =\operatorname{Hom}_S(I_\epsilon,S)=S,
 \qquad M_\epsilon^{**}=S,
 \qquad
 \operatorname{im}(M_\epsilon\longrightarrow M_\epsilon^{**})
        =I_\epsilon.
 \tag{SC.11}
\]

The evaluation map is exactly the extension quotient followed by
\(I_\varepsilon\to S\). These identifications reduce to the corresponding
identifications for \(M\), because their generators are those actual
hull maps. They hold over the nonreduced ring \(S\); regularity of \(S\)
has not been asserted or used.

\subsection{4. The tangent map and its exact kernel}\label{SF-E847:s4}

The \(D\)-flat family defines a \(\mathbf C\)-linear map

\[
 L:A^4\oplus A^6\longrightarrow\operatorname{Ext}^1_R(M,M).
 \tag{SC.12}
\]

One elementary description takes the self-extension
\(0\to M\to M_\varepsilon\to M\to0\), with inclusion induced by \(\varepsilon\). Equivalently,
differentiate the just constructed free differential: its first-order
coefficient is a degree-one endomorphism cocycle of the fixed resolution.
It is linear in \(\alpha\) and \(\beta\), and boundaries are exactly the
first-order changes of trivialization. This proves linearity of \(L\).

For completeness, a zero self-extension gives exactly a trivial \(D\)-flat
module deformation, with identity special fiber. If \(s:M\to M_\varepsilon\) is
an \(R\)-linear splitting, the map
\((m_0+\varepsilon m_1)\mapsto s(m_0)+\varepsilon s(m_1)\) is a \(D\)-linear isomorphism
from \(M\otimes D\) to \(M_\varepsilon\), reducing to the identity. Its inverse
also gives the converse. This checks the stated equivalence directly.

Suppose \(L(\alpha,\beta)=0\). The resulting \(S\)-linear isomorphism
\(M_\varepsilon\simeq M\otimes D\), with identity special fiber, induces a unit
automorphism of the relative double dual \(S\). By (SC.11) it carries
\(I_\varepsilon\) to \(I\otimes D\). A unit times an ideal is the same ideal,
so \(I_\varepsilon=I\otimes D\). In particular \(\varepsilon\alpha_i\) belongs to
\(I\otimes D\) for every \(i\), and therefore \(\alpha_i=0\) in \(A\). Thus \(\alpha=0\).
For our fixed linear lifts, the ideal family is now exactly constant.

The induced automorphisms on the torsion \(A\otimes D\) and on the hull \(S\)
are multiplication by units \(1+\varepsilon t\) and \(1+\varepsilon r\) respectively.
The map on the hull-image ideal is the restriction of the latter.
Equality of the two pushed/pulled extension classes is therefore

\[
 (1+\epsilon t)(b+\epsilon\beta)
       =(1+\epsilon\bar r)b,
 \qquad\hbox{so}\qquad
 \beta=(\bar r-t)b.
 \tag{SC.13}
\]

Conversely, if \(\beta=sb\), multiplication by \(1-\varepsilon s\) on the torsion
identifies the extension \(b+\varepsilon\beta\) with \(b\), with identity quotient
and identity special fiber. This can be read either from the defining
equivalence of Yoneda extensions or from the triangular resolution,
multiplying its top Koszul summand by \(1-\varepsilon s\). Hence

\[
 \ker L=\{0\}\oplus Ab,
 \qquad \dim\ker L=\ell-a,
 \qquad e_1\ge10\ell-(\ell-a)=9\ell+a.
 \tag{SC.14}
\]

Only this explicit family of deformations is counted. No claim that all
deformations of \(M\) preserve the scalar-CI form is needed.

\subsection{5. The finite-supported hull-cone identity}\label{SF-E847:s5}

Put \(C=\operatorname{cone}(\phi)\) in the local distinguished triangle

\[
 M\xrightarrow{\phi}R\longrightarrow C\longrightarrow M[1].
 \tag{SC.15}
\]

Its cohomology is \(A\) in degree \(-1\) and \(A\) in degree zero, and hence its
finite-length Euler characteristic is zero. Both \(M\) and \(C\) are perfect;
the explicit triangular Koszul resolution of \(M\) has length four.
Set \(F=R\operatorname{Hom}_R(M,C)\). This complex has finite-length cohomology, supported
in degrees \(-1\) through \(4\). The finite filtration obtained by a free
resolution of \(M\) proves

\[
 \chi_{\mathrm{length}}(F)
   =\left(\sum_i(-1)^i\operatorname{rank}L_i\right)
             \chi_{\mathrm{length}}(C)=0.
 \tag{SC.16}
\]

This is an Euler characteristic of finite-supported complexes only.
It does not cancel infinite dimensions in \(\operatorname{End}_R(M)\).

Write \(h_j=\dim\operatorname{Ext}^j_R(M,C)\). The two hull short exact sequences and
\(\operatorname{Ext}^j_R(A,R)=0\) for \(j<4\) give \(\operatorname{Ext}^1_R(M,R)=\operatorname{Ext}^2_R(M,R)=0\).
Apply \(\operatorname{Hom}_R(M,-)\) to (SC.15), using the surjection \(\operatorname{End}_R(M)\to R\)
from (SC.7). The first terms give

\[
 h_{-1}=k,\qquad h_0=e_1,\qquad h_1=e_2.
 \tag{SC.17}
\]

To compute the remaining terms, use local duality on \(F\), not on the
non-finite-supported endomorphism complex. After a choice of coordinate
volume, finite-support local duality and perfect tensor-\(\operatorname{Hom}\) adjunction
give

\[
 \operatorname{Ext}^j_R(M,C)^\vee
      \simeq\operatorname{Ext}^{4-j}_R(C,M).
 \tag{SC.18}
\]

Here the dual is the finite-dimensional complex-vector-space dual.
More explicitly, local duality identifies the dual of the cohomology
of \(F\) with that of \(R\operatorname{Hom}_R(F,R[4])\). Perfectness of \(M\) identifies the latter
complex with \(R\operatorname{Hom}_R(C,M)[4]\). Derived completion causes no change:
all the cohomology in question has finite length and \(R\) is complete.
The precise general input is Stacks, Section 47.18, tag 0A81,
Theorem 47.18.3 and Lemma 47.18.4,\cite{SFStacksE847LocalDuality} with their complete proofs freshly
read in this run. For \(R=\mathbf C[[x]]\), its normalized dualizing complex is
\(R[4]\); the injective-hull dual of a finite-length module is its \(\mathbf C\)-dual
via the coordinate residue, as in Lipman Lemma 7.3/Theorem 7.4,
pp.~64--69.\cite{SFLipman1984} No local cohomology of \(\operatorname{End}_R(M)\) is discarded.

Apply \(\operatorname{Hom}_R(-,M)\) to (SC.15). The map \(M=\operatorname{Hom}_R(R,M)\to\operatorname{End}_R(M)\) sends
\(m\) to the endomorphism \(z\mapsto\phi(z)m\). Its kernel is exactly \(A\). Indeed
every element of the torsion \(A\) is annihilated by \(I\). Conversely, if
\(Im=0\), applying \(\phi\) gives \(I\phi(m)=0\) in the domain \(R\). Since \(I\) is
nonzero, \(\phi(m)=0\) and \(m\) belongs to \(A\). Thus \(\operatorname{Hom}_R(C,M)=A\).
The image of \(M\) in \(\operatorname{End}_R(M)\) is isomorphic to \(I\); under the projection
\(\operatorname{End}_R(M)\to R\) its image maps isomorphically onto \(I\). Consequently

\[
 \begin{split}
 &\dim\operatorname{Hom}_R(C,M)=\ell,\\
 &0\longrightarrow K\longrightarrow\operatorname{Ext}^1_R(C,M)
                    \longrightarrow A\longrightarrow0,\\
 &\operatorname{Ext}^2_R(C,M)\simeq\operatorname{Ext}^1_R(M,M).
 \end{split}
 \tag{SC.19}
\]

The last isomorphism uses \(\operatorname{Ext}^i_R(R,M)=0\) for \(i>0\). Formula (SC.18)
now gives \(h_4=\ell\), \(h_3=k+\ell\) and \(h_2=e_1\). Combining all six finite
cohomology groups with (SC.16) yields

\[
 0=-k+e_1-e_2+e_1-(k+\ell)+\ell
    =2e_1-e_2-2k.
 \tag{SC.20}
\]

This proves (SC.4). One may also deduce \(e_3=e_1\) from the same supported
duality and the long exact sequences, but that identity is not used as
an unproved self-duality premise.

\subsection{6. A canonical global injection and its six-dimensional trace image}\label{SF-E847:s6}

Let \(E\) and \(H\) be as in the global theorem, and let \(\mathcal C=\operatorname{cone}(E\to H)\).
All its cohomology is supported on a finite set \(Z\). Since \(E\) is perfect,
the complex \(R\operatorname{Hom}(E,\mathcal C)\) has finite-supported coherent cohomology.
Its support spectral sequence has no positive sheaf-cohomology rows.
Restriction and completion therefore give natural isomorphisms

\[
 \operatorname{Ext}^j_X(E,\mathcal C)
    \simeq\bigoplus_{p\in Z}
       \operatorname{Ext}^j_{\widehat{\mathcal O}_{X,p}}(M_p,C_p).
 \tag{SC.21}
\]

At every support point, the free hull and finite kernel/cokernel give
\(\operatorname{Ext}^1_R(M_p,R^r)=\operatorname{Ext}^2_R(M_p,R^r)=0\). The local hull triangle identifies
\(\operatorname{Ext}^1_R(M_p,C_p)\) with \(\operatorname{Ext}^2_R(M_p,M_p)\). The global preceding arrow
\(\operatorname{Ext}^1_X(E,H)\to\operatorname{Ext}^1_X(E,\mathcal C)\) is actually zero: in its commutative
restriction/completion square the bottom source is zero, and the
right vertical map is the isomorphism (SC.21). Thus the global
connecting map gives a natural injection

\[
 j:\bigoplus_{p\in Z}\operatorname{Ext}^2_R(M_p,M_p)
       \hookrightarrow\operatorname{Ext}^2_X(E,E).
 \tag{SC.22}
\]

It is a section of restriction to the local \(\operatorname{Ext}\) groups and is natural
for every isomorphism of the actual hull triangle. This repeats the
full E846 injection argument to specify the map being used.

Let \(u\) be any class in the source, and let \(a\in\operatorname{Ext}^1_X(E,\mathcal C)\)
be its corresponding preimage before the connecting map. For a global
Hochschild degree-two class \(\gamma\), associativity and naturality of
composition give, up to the harmless fixed cyclic sign,

\[
 \operatorname{Tr}_E\bigl(j(u)\rho_E(\gamma)\bigr)
   =\operatorname{Tr}_E\bigl(\delta\,[a\rho_E(\gamma)]\bigr).
 \tag{SC.23}
\]

The bracket belongs to \(\operatorname{Ext}^3_X(E,\mathcal C)\), so (SC.21) makes it depend
only on the actual completed restrictions of \(\rho_E(\gamma)\). The HKR
components \(H^1(T_X)\) and \(H^2(\mathcal O_X)\) restrict to zero on the affine formal
local schemes. Naturality of the Hochschild characteristic map gives
the same zero restriction for their evaluations. The factorization
through (SC.21) ensures that no hidden Cech boundary contributes to
(SC.23). Thus the pairing factors through

\[
 HH^2(X)\longrightarrow H^0(X,\Lambda^2T_X),
 \qquad \dim H^0(X,\Lambda^2T_X)=\binom42=6.
 \tag{SC.24}
\]

On an abelian fourfold the supplied CY4 Serre-trace/evaluation adjunction
identifies the annihilator of all these evaluations with the kernel of
the full semiregularity map. Equivalently, its nondegenerate Hochschild
pairing transposes \(\rho^2\) and \(\sigma^2\). The exact input is DEG\_L S20,
first proof paragraph, and S21.6 through the adjunction immediately
before S21.8, already read in the E846 audit. Stacks 48.27 Lemma 48.27.1(5),(7),
Remarks 48.27.2--3 and Lemma 48.27.4\cite{SFStacksE846Duality} give the ordinary perfect-complex
Serre pairing. This argument uses the stronger adjunction explicitly;
the first-order formula \(\sigma\rho=\operatorname{contr}(\operatorname{ch})\) alone is insufficient.

It follows that on any \(d\)-dimensional subspace of the source of \(j\),
the full semiregularity kernel has dimension at least \(d-6\). No internal
weight, coordinate homogeneity, or residue-weight cancellation is
being used in this rank bound.

\subsection{7. Specified finite actions and invariant multiplicity spaces}\label{SF-E847:s7}

We now prove that a subspace of dimension \(e_2\) from (SC.4) can be chosen
inside invariant local \(\operatorname{Ext}\), without assuming a grading-preserving action.
Consider only the point stabilizer \(G_p\) of the specified translation/tensor
action. A translation fixing \(p\) is the identity translation. Therefore
the base action at this stage is \(R\)-linear. The Picard factors are torsion
line bundles. In flat local frames their coherent multiplication
isomorphisms differ by constants, giving one constant projective cocycle
\(\alpha\) in \(Z^2(G_p,\mathbf C^*)\). The resulting action on \(\operatorname{Ext}\) is honest.

The twisted finite group algebra \(\mathbf C_\alpha[G_p]\) is semisimple: average
conjugates of a linear projection; the projective scalar factors cancel.
Its simple representations \(S_\lambda\) yield the actual \(R\)-linear decomposition
\(M_p=\bigoplus_\lambda S_\lambda\otimes N_\lambda\), with action on \(S_\lambda\) alone.
The modules \(N_\lambda\) are finite over \(R\) and are direct summands of finite
direct sums of \(M_p\).

The complete-local Krull--Schmidt fact used in E846 applies here unchanged.
For specificity, endomorphism rings of finite \(R\)-modules are finite and
\(\mathfrak m\)-adically complete; idempotents lift from their finite-dimensional
reductions and then through all powers of \(\mathfrak m\). This supplies decompositions
into summands with local endomorphism rings. Uniqueness follows by
expressing the identity on an indecomposable summand as a sum of composites
through another decomposition: a sum of nonunits cannot be the identity
in a local ring, so one composite is invertible and its summand can be
split off. In the present case \(\operatorname{End}_R(M_b)\) and \(\operatorname{End}_R(R)\) are local by
Section 2, and \(M_b\) is not \(R\). Consequently every \(N_\lambda\) is a sum of
copies of \(M_b\) and \(R\). Regrouping the actual isotypic decomposition gives

\[
 M_p\simeq(M_b\otimes V)\oplus(R\otimes W),
 \qquad \dim V=m>0,
 \tag{SC.25}
\]

with projective multiplicity representations having the same scalar
cocycle. No original choice of local splitting or resolution is asserted
to be invariant. The characteristic and hull maps are transported through
this actual isomorphism.

For every \(u\) in \(\operatorname{Ext}^2_R(M_b,M_b)\), the class \(u\otimes\operatorname{id}_V\), zero on
the free summand, is invariant and nonzero when \(u\) is nonzero. This gives
an \(e_2\)-dimensional invariant subspace. Transport over the orbit of \(p\)
and sum its components. Stabilizer invariance makes this independent
of coset representatives; the resulting map remains injective because
its component at \(p\) is the original class. The injection (SC.22) is
equivariant by naturality of the canonical hull triangle.

By (SC.24), the ordinary trace pairing on this \(e_2\)-dimensional invariant
subspace has rank at most six, including after summing all orbit components.
The higher HKR components vanish at every support point. In the supplied
effective translation/tensor framework, S21 identifies this ordinary
semiregularity map on invariant \(\operatorname{Ext}\) with the identity-sector map; the
other sectors vanish as proved there. Thus its kernel is a kernel for
the weak \(G\)-semiregularity map itself. We obtain the bound

\[
 \dim\ker(\sigma^2_E|_{\operatorname{Ext}^2(E,E)^G})
       \ge e_2-6\ge10\ell-6>0.
 \tag{SC.26}
\]

This proves the global theorem. Arbitrary finite automorphisms with
nontrivial derivative at a fixed point are outside this action scope.

\subsection{8. Actual coherent balanced objects for every stated scalar type}\label{SF-E847:s8}

These local types are algebraic at any chosen point of an abelian
fourfold. Because \(I\) has finite colength, it contains \(\mathfrak m^N\) for some \(N\).
The quotient by \(I\) is therefore already a quotient of the algebraic
local ring modulo \(\mathfrak m^N\), using its canonical isomorphism with the same
finite jet of the completion. It defines a closed finite subscheme \(Z\).
The finite-support spectral sequence for \(R\operatorname{Hom}(\mathcal I_Z,\mathcal O_Z)\) identifies its
global \(\operatorname{Ext}^1\) with the completed local \(\operatorname{Ext}^1\). Hence every \(b\) in \(A^6\) gives
an actual extension \(0\to\mathcal O_Z\to M_Z\to\mathcal I_Z\to0\), with completed module \(M_b\),
free hull \(\mathcal O_X\), and equal-length torsion and hull quotient. Its full
Chern character is 1 by additivity.

For an actual indecomposable balanced example, take the previously proved
rank 32 hull \(H\) on \(E_i^4\) from E815/E820, with \(\operatorname{ch}(H)=32v_{1/2}\). Its seven
line types have multiplicities \(24,2,2,1,1,1,1\). Distinct types have no \(\operatorname{Hom}\)
by the primitive index-two calculation in E813/E815. Thus \(\operatorname{End}(H)\) is
the constant block algebra \(\operatorname{Mat}_{24}(\mathbf C)\) \(\oplus\) \(\operatorname{Mat}_2(\mathbf C)\) \(\oplus\) \(\operatorname{Mat}_2(\mathbf C)\)
\(\oplus\) \(\mathbf C^4\).

At 33 distinct points choose algebraic finite subschemes with the given
completed ideal \(I\). In local frames of the 32 line summands, map \(H\) onto
one copy of \(A\) at each point, using the 32 coordinate rows and then the
all-ones row. These are actual surjective morphisms to the finite sheaf
\(Q\). Let \(K_{\mathrm{global}}\) be the kernel, and set \(T=Q\). At each point \(K_{\mathrm{global}}\) has
completed module \(I\oplus R^{31}\). Choose the scalar class \(b\) on its \(I\)
factor; the free factors have no local \(\operatorname{Ext}^1\) to \(A\). The same support
spectral sequence supplies an actual global extension
\(0\to T\to E\to K_{\mathrm{global}}\to0\). Its local module is \(M_b\oplus R^{31}\), its hull
is \(H\), and its torsion and hull quotient are \(T\) and \(Q\). Thus \(\operatorname{ch}(E)=\operatorname{ch}(H)\).

Every idempotent endomorphism of \(E\) induces a constant block-matrix
endomorphism of \(H\) preserving these quotient rows. The 32 coordinate
rows force it to be diagonal, and the all-ones row forces it to be
\(\lambda\) times the identity, with \(\lambda\) equal to \(0\) or \(1\). If \(\lambda=0\),
the torsion idempotent at each point is multiplication by \(t\) in \(A\) with
\(tb=0\). Since \(b\) is nonzero, \(t\) has zero residue. The only such Artin-local
idempotent is zero. The original idempotent then has image in \(T\) and
vanishes on \(T\), so its square is zero and it is zero. If \(\lambda=1\),
apply the same argument to its complement. Therefore \(E\) is indecomposable.
It is coherent and hence ordinary gluable. The identity group is an
actual action; no additional finite action is assumed to exist.

Theorem E847b excludes weak semiregularity for the identity action and
for every further specified action if it exists. These actual examples
do not, by themselves, prove strict separation from the earlier Artin
obstructions. That comparison remains a separate gate.

\subsection{9. Source, verification and boundary record}\label{SF-E847:s9}

The new local proof is the explicit ten-parameter-vector Koszul family,
its exact tangent kernel, and the finite-cone Euler identity. It does not
assume a moduli stack dimension, unobstructedness of all local module
deformations, or local self-duality of an endomorphism complex with
rank-one cohomology. Every use of local duality is on the finite-supported
complex \(F\).

The ordinary global injection, support factorization and actual projective
stabilizer decomposition have been spelled out, with their previously
checked E846 and supplied S20/S21 interfaces identified. The general
duality source passages stated in Section 5 were freshly reopened and
read in full, as were the invoked perfect-complex pairing passages in
Stacks 48.27. No new full reading of Lipman's monograph, Buchweitz--Flenner,
or the entire supplied input is claimed for this proof.

There is no numerical certificate claimed for the universal dimension
inequality at this draft stage. Its proof is algebraic. Matrix-valued
extension classes, local types not actually isomorphic to (SC.5),
non-locally-free hulls, positive-dimensional defects, and general perfect
complexes are outside the theorem. The full proof has passed root and independent review. No independent
main-gate movement is claimed without a concrete strict comparison.
